\section{综合讨论}

\subsection{三任务横向对比}

三项任务共享同一数据集、同一预处理流水线与同一CatBoost主模型，但预测难度呈现出三种截然不同的格局。特征信号在不同预测目标上的异质性分布构成这一差异的根源。表\ref{tab:overall}汇总了关键指标。

\begin{table}[H]
\centering
\caption{三项任务综合对比}
\label{tab:overall}
\begin{tabular}{lcccl}
\toprule
\textbf{维度} & \textbf{任务一} & \textbf{任务二} & \textbf{任务三} & \textbf{说明} \\
\midrule
预测类型 & 三分类 & 二分类 & 二分类 & \\
验证集ACC & 0.8633 & 0.6498 & 0.8550 & \\
超过多数类基线 & +27.0pp & +0pp & +34.7pp & 信号强度对比 \\
ROC-AUC & -- & $\approx$0.50 & 0.9042 & 排序能力对比 \\
CV折间方差 & 0.0025 & 0.0000 & 0.0014--0.0026 & 信号稳定性 \\
特征信号强弱 & 强 & \textbf{无} & 强 & \\
\bottomrule
\end{tabular}
\end{table}

任务一与任务三的特征信号较强，CatBoost在交叉验证和独立验证集上均稳定超越多数类基线。任务三同时达到较高的AUC（0.9042）和ACC（0.8550），边界分析中出现的V形错误曲线为这一表现提供了诊断性确认——模型受真实信号驱动，而非噪声拟合。任务二呈现完全不同的图景：四种候选模型的ROC-AUC全部在0.50附近，交叉验证折间方差严格为零。在现有特征集合范围内，家庭充电安装与否无法被有效预测。

\subsection{CatBoost统一技术栈的优势}

三项任务均采用CatBoost作为主模型，这一选择服务于以下四个方面的考虑。

\begin{enumerate}[itemsep=0pt]
\item \textbf{方法学一致性}：统一主模型使跨任务比较不受模型选型差异的干扰，保证了论文分析逻辑的内在连贯。
\item \textbf{实现效率}：分类特征直接在CatBoost内部声明，避免独热编码带来的维度膨胀；三项任务共享common.py、features.py等公共模块，代码复用程度较高。
\item \textbf{诊断价值}：相同模型能力上限意味着跨任务表现差异只能归因于特征来源的不同——"模型选择不当"这一替代解释得以排除。
\item \textbf{可解释性}：CatBoost内置的PredictionValuesChange特征重要性度量和SHAP值为全部三个任务提供了统一的解释框架，便于论文中并行展开分析。
\end{enumerate}

\subsection{特征工程策略反思}

三项任务中构造了大量业务交互特征，其效果因任务信号基础不同而两极分化。

\begin{itemize}[itemsep=0pt]
\item \textbf{任务一}：knowledge\_anxiety\_gap与green\_tech\_score两个业务交互特征占据特征重要性前两位（合计41.7分），表明认知维度上的交互构造具备较强的分类能力。
\item \textbf{任务二}：充电稀缺度、城市交互项等业务特征对预测近乎无贡献。底层单变量与目标间的相关性已接近零，业务特征因此失去了构造基础。
\item \textbf{任务三}：long\_distance\_ratio、knowledge\_buffer等业务特征出现在Top 15中，但位次居于中后段。CatBoost的对称树结构已在一定程度上内化了这些交互效应，业务特征在此场景下的作用更偏向显式化已知交互关系。
\end{itemize}

三项任务的横向对照指向一个共有的约束条件：特征工程的收益上限受制于底层信号的存在性。原始特征与目标无关时（任务二），任何交互构造均难以产生区分力。而原始特征信号足够时（任务三），CatBoost自身的树结构已可有效捕捉非线性交互，业务特征提供的是边际增益。

\subsection{任务二负结果的学术价值}

任务二是三项任务中唯一未能突破基线的子问题，但该负结果本身具备不可忽视的方法论价值。相当数量的数据科学竞赛报告倾向于仅呈现成功案例；本文选择完整记录从EDV预诊断、三方案消融到最终判定的全流程，目的如下：

\begin{enumerate}[itemsep=0pt]
\item 展示特征来源信号不足时的诊断与排除流程——先通过多模型比照确认信号是否真实存在，再决定是否投入超参数搜索，避免因果倒置。
\item 区分"模型选择不合理"与"数据本身缺乏信息"两种性质截然不同的困境，证明系统方法论可以对二者作出可靠区分。
\item 为后续研究提供一份可复用的特征有效性预诊断模板：EDV先行，调参在后。
\end{enumerate}
